\documentclass[12pt]{article}
\usepackage[utf8]{inputenc}
\usepackage[T1]{fontenc}
\usepackage{hyperref}
\title{Toram Language Resources for K-12 Curriculum Development}
\author{Compiled by Claude Sonnet 5 (claude-sonnet-5) for Andr\'as Kornai}
\date{}
\begin{document}
\maketitle

\section{Introduction}
Toram (also Torom, Torum; ISO 639-3 \texttt{trj}) is an East Chadic language of the Afro-Asiatic family, spoken in the Batha region of central Chad near Lake Fitri. Speakers have shifted toward Chadian Arabic, and younger people are reported to speak Arabic. Speaker counts differ widely between sources (under 4,000; 8,500; 28,000 for the ethnic group). Little material exists.

\section{Primary Educational Resources}
None were found in the search results.

\section{Linguistic and Reference Materials}
\begin{itemize}
\item Two papers on Toram and related languages appear in the Linguistica Posnaniensia series, found through the search. I could not confirm their titles or authors from the results. First paper: \url{https://pressto.amu.edu.pl/index.php/linpo/article/download/36985/31735/79722}. Second paper: \url{https://pressto.amu.edu.pl/index.php/linpo/article/download/39097/33249/86795}
\item General overview: \url{https://en.wikipedia.org/wiki/Toram_language}
\item Endangered Languages Project entry: \url{https://www.endangeredlanguages.com/lang/582?hl=de}
\end{itemize}

\section{Religious and Community Resources}
No scripture or community organisation was found. Joshua Project lists a people-group profile: \url{https://joshuaproject.net/people_groups/15567}

\section{Digital and Audio Resources}
None were found in the search results.

\section{Curriculum Development Context}
No Toram teaching materials, orthography or audio were found, and the language is endangered. Any curriculum would require new documentation work with speakers, and the scholarly papers above may help with the linguistic groundwork.
\end{document}